% Chapter: Genesis: The Edge of Chaos (1988--1992)
% Track: 1
% Plan: 0002

\settrack{trackone}

\chapter{Genesis: The Edge of Chaos (1988--1992)}
\label{t1:ch01:genesis}
\label{pos13:genesis}
\aidraftchapter

% Abstract epigraph
\begin{quote}\small\textit{%
``We propose that non-abelian anyonic quasiparticles in a fractional quantum Hall system, subjected to structured electromagnetic perturbation, can undergo autocatalytic self-organization at the edge of chaos, producing an emergent topological quantum neural network\ldots''
}\end{quote}

\vspace{0.5cm}

This chapter traces the emergence of self-organized quantum computation from first principles: how autocatalytic chemical networks, constrained by topology, produce structures capable of information processing. The argument proceeds from Kauffman's autocatalytic sets through the quantum Hall effect to the threshold where chemistry becomes computation.

Hasslacher at Los Alamos. Soliton dynamics on lattice-gas automata running on Hillis's Connection Machine. Freedman's 1988 independent insight that topology could protect quantum information. Kauffman's autocatalytic sets. Wolfram's computational universality. Five minds, five disciplines, one synthesis that none of them could have reached alone.

\chapterreturn
